% Partie : histoire
% Chapitre : biographies
% Section : Bertrand Russell

\documentclass[../../../include/open-logic-section]{subfiles}

\begin{document}

\olfileid[fr]{his}{bio}{rus}

\olsection{Bertrand Russell}

\olphoto{russell-bertrand}{Bertrand Russell}

Bertrand Russell est considéré comme l'un des fondateurs de la
philosophie analytique moderne. Né le 18 mai 1872, il est connu pour
ses travaux en philosophie et en logique, mais aussi pour ses nombreux
ouvrages de vulgarisation sur des sujets variés. Il fut également un
militant politique passionné tout au long de sa vie.

Russell naquit à Trellech, dans le Monmouthshire, au pays de Galles.
Ses parents appartenaient à la noblesse britannique. Libres-penseurs,
ils s'étaient même liés d'amitié avec les radicaux de Boston de leur
époque. Malheureusement, ils moururent alors que Russell était encore
jeune, et celui-ci fut confié à ses grands-parents. Il y reçut une
éducation religieuse, ce que ses parents avaient voulu éviter à tout
prix. Sa grand-mère était très stricte en matière de morale. Pendant
son adolescence, il étudia surtout à domicile avec des précepteurs.

L'influence de Russell sur la philosophie analytique, et
particulièrement sur la logique, est immense. Il étudia les
mathématiques et la philosophie au Trinity College de Cambridge, où
le mathématicien et philosophe Alfred North Whitehead exerça une
influence sur lui. En 1910, Russell et Whitehead publièrent le premier
volume des \emph{Principia Mathematica}, dans lequel ils défendaient
la thèse selon laquelle les mathématiques peuvent être ramenées à la
logique. Russell publia ensuite des centaines de livres, d'essais et
de brochures politiques. En 1950, il reçut le prix Nobel de
littérature.

Russell s'engagea profondément dans la politique et l'action sociale.
Pendant la Première Guerre mondiale, ses activités pacifistes et ses
protestations lui valurent une arrestation et une condamnation
à six mois de prison. Il fut libéré en septembre 1918, six semaines
avant la fin de sa peine. En prison, il put lire et écrire, et déclara
avoir trouvé cette expérience
«~assez agréable~». Il poursuivit son engagement pour la paix et fut de nouveau
incarcéré en 1961 pour avoir participé à une manifestation en faveur
du désarmement nucléaire. En 1948, il avait aussi survécu à un
accident d'avion dont les seuls rescapés se trouvaient dans la section
fumeurs. Russell affirma donc devoir la vie au tabac. Il se maria
quatre fois, mais avait aussi la réputation d'entretenir des liaisons
extraconjugales. Il mourut le 2 février 1970, à 97 ans, à
Penrhyndeudraeth, au pays de Galles.

\begin{reading}
Russell écrivit une autobiographie en trois volumes, couvrant la
période 1872--1967 \citep{Russell1967,Russell1968,Russell1969}.
Le Bertrand Russell Research Centre de l'université McMaster conserve
les archives de Russell. Son site, cité dans \citet{Duncan2015},
présente les volumes des \emph{Collected Papers}, avec des index
interrogeables, ainsi que les projets d'archivage. L'article de
Russell \emph{On Denoting} \citep{Russell1905} est un classique de la
philosophie analytique du XXe siècle.

La notice consacrée à Russell dans la Stanford Encyclopedia of
Philosophy \citep{Irvine2015} contient des extraits sonores où il
parle du désir et de la théorie politique. De nombreux entretiens
filmés avec Russell sont accessibles en ligne. Pour l'entendre parler
du tabac et de l'accident d'avion, voir par exemple \citet{RussellND}.
Certaines de ses œuvres, dont l'\emph{Introduction to Mathematical
Philosophy}, sont disponibles gratuitement en livres audio sur
\citet{LibriVoxND}.
\end{reading}

\end{document}
